\documentclass{beamer}
\usepackage[nju,zh]{collegeBeamer}

% Chinese needs XeLaTeX (or LuaLaTeX).  The package picks a CJK font
% automatically -- Noto Sans/Serif CJK SC, Source Han Sans/Serif SC or
% Microsoft YaHei -- but you can override it:
%   \setCJKmainfont{Source Han Serif SC}
%   \setCJKsansfont{Source Han Sans SC}

% ---------------------------------------------------------------------------
%  Meta data -- replace the placeholders below with your own.
% ---------------------------------------------------------------------------
\title{学院 Beamer 模板}
\subtitle{标题、副标题与作者信息的位置示意}
\author{作者姓名\quad 合作者姓名\quad 课题组}
\date{2026年10月12日}

% ---------------------------------------------------------------------------
%  Document body.  Every page below is a demonstration of one element of the
%  template -- the wording is deliberately generic so that you can overwrite
%  it without touching the layout.
% ---------------------------------------------------------------------------
\begin{document}

\maketitle

\begin{frame}{模板说明}
    本模板基于 \href{https://www.overleaf.com/latex/templates/sintef-presentation/jhbhdffczpnx}{SINTEF Presentation}
    模板重制，视觉风格参考南京大学版演示稿：整版主色标题页、页眉处的章节序号、
    主色分隔线、浅色信息面板与卡片，以及右下角的页码。

    \vspace{0.3cm}
    所有学校通用，只需更换包选项即可切换配色与校徽：
    \texttt{nju}、\texttt{polyu}、\texttt{szu}、\texttt{swu}、\texttt{hit}、\texttt{saes}、\texttt{zafu}、\texttt{red}。
\end{frame}

\section{页面元素}

\begin{frame}{页眉与项目符号}{页眉显示当前章节，帧标题作为引导句}
    项目符号列表的示例。每一行都是一段说明文字，用来占位并展示行距、
    缩进与强调色的效果：

    \begin{itemize}
        \item 第一级要点：普通正文使用衬线字体，页眉与标题使用无衬线字体
        \item 第二级要点：数学公式、表格与代码块都能直接排版
        \item 第三级要点：需要强调时使用 \ctextbf{主色加粗} 或 \cemph{主色斜体}
    \end{itemize}

    \vspace{0.35cm}
    \keypoint{关键句：空心箭头加粗体主色，用于强调一段结论。}
\end{frame}

\begin{frame}{卡片与信息面板}
    \begin{columns}[T,onlytextwidth]
        \begin{column}{0.23\textwidth}
            \card{01}{步骤一}{卡片标题与正文，卡片用于并列展示若干个并列要点。}
        \end{column}
        \begin{column}{0.23\textwidth}
            \card{02}{步骤二}{卡片标题与正文，卡片用于并列展示若干个并列要点。}
        \end{column}
        \begin{column}{0.23\textwidth}
            \card{03}{步骤三}{卡片标题与正文，卡片用于并列展示若干个并列要点。}
        \end{column}
        \begin{column}{0.23\textwidth}
            \card{04}{步骤四}{卡片标题与正文，卡片用于并列展示若干个并列要点。}
        \end{column}
    \end{columns}

    \vspace{0.4cm}
    \callout[说明]{信息面板：浅色底加左侧主色竖条，用来放一段需要突出的说明文字。}
\end{frame}

\begin{frame}{表格与数据卡}
    {\sffamily\bfseries\color{maincolor}表头使用主色加粗，正文行可加浅色斑马纹。}
    \vspace{0.25cm}

    \zebra
    \setlength{\tabcolsep}{13.5pt}
    \begin{tabular}{lrrrrrr}
        \rowcolor{mainlight}
        \thead{第一列} & \thead{第二列} & \thead{第三列} & \thead{第四列} &
        \thead{第五列} & \thead{第六列} & \thead{第七列} \\
        一 & 1.00 & 2.00 & 3.00 & 4.00 & 5.00 & 6.00 \\
        二 & 1.50 & 2.50 & 3.50 & 4.50 & 5.50 & 6.50 \\
        三 & 2.00 & 3.00 & 4.00 & 5.00 & 6.00 & 7.00 \\
        四 & 2.50 & 3.50 & 4.50 & 5.50 & 6.50 & 7.50 \\
        五 & 3.00 & 4.00 & 5.00 & 6.00 & 7.00 & 8.00 \\
    \end{tabular}

    \vspace{0.35cm}
    \begin{columns}[T,onlytextwidth]
        \begin{column}{0.30\textwidth}
            \stat{第一个数据卡}{标签 + 数值}
        \end{column}
        \begin{column}{0.30\textwidth}
            \stat{第二个数据卡}{标签 + 数值}
        \end{column}
        \begin{column}{0.30\textwidth}
            \stat{第三个数据卡}{标签 + 数值}
        \end{column}
    \end{columns}
\end{frame}

\section{公式与代码}

\begin{frame}{数学公式}
    公式页的示例。行间公式与行内公式都可以直接使用，数学字体由模板统一设置。

    \vspace{0.3cm}
    \[
        S_t = \{\, i : i \text{ 在时间窗内出现过条件}\,\},
        \qquad
        R_t = \frac{1}{|S_t|}\sum_{i \in S_t} r_{i,t}
    \]
    \[
        I_t = I_{t-1}\,\bigl(1 - R_t\bigr)
    \]

    \vspace{0.3cm}
    \callout[口径说明]{信息面板同样可以放在公式下方，用来解释符号的含义与口径。}
\end{frame}

\begin{frame}[fragile]{代码块}
\begin{codepanel}{示例代码}
\begin{lstlisting}[language=Python]
def summarise(df):
    """一行摘要文字，用来占位。"""
    grouped = df.groupby(df.index.dayofweek)
    return grouped.agg(n=("value", "size"),
                       mean=("value", "mean"))
\end{lstlisting}
\end{codepanel}
\end{frame}

\section{自定义组件}

\subsection{面板与引导句}

\begin{frame}{双栏面板}
    \begin{columns}[T,onlytextwidth]
        \begin{column}{0.62\textwidth}
            \callout[左侧面板]{左栏的面板与右栏的面板顶部对齐，高度随内容伸缩。}
        \end{column}
        \begin{column}{0.34\textwidth}
            \callout[右侧面板]{右栏的面板与左栏的面板顶部对齐，高度随内容伸缩。}
        \end{column}
    \end{columns}
\end{frame}

\begin{frame}{结论页}
    \keypoint{引导句：用于在一页的开头点明最重要的那条结论。}

    \vspace{0.3cm}
    \begin{itemize}
        \item 要点一：一条普通正文内容
        \item 要点二：一条普通正文内容
        \item 要点三：一条普通正文内容
    \end{itemize}

    \vspace{0.35cm}
    \callout[下一步]{结尾的信息面板：写下一步计划、致谢或联系方式。}
\end{frame}

\section{参考文献}

\begin{frame}{References}
    \begin{enumerate}
        \item Author A, Author B. Title of the first reference[J]. Journal Name, 2020, 12(3): 45--67.
        \item Author C, Author D. Title of the second reference[C]. Proceedings of the Conference, 2021: 101--110.
        \item Author E. Title of the third reference[M]. City: Publisher, 2022.
    \end{enumerate}
\end{frame}

\QApage

\end{document}
